\section{Class Diagram và Method Descriptions}

Phần này trình bày lớp và method dùng để hiện thực toàn bộ phạm vi. Class diagram không chỉ là danh sách entity; cần có boundary/application service/domain service/repository/adapter khi các trách nhiệm đó tồn tại trong thiết kế. Mọi method xuất hiện trên diagram phải có mô tả tương ứng.

\subsection{Class diagram tổng thể}

\ArtifactBrief{SYSTEM CLASS DIAGRAM}{CD-SEMH-OVERVIEW}{Thể hiện quan hệ giữa các package/phân hệ, shared types, adapter IoT/map và dependency direction. Không nhồi toàn bộ thuộc tính nhỏ vào một hình duy nhất.}{Class diagram tổng thể của SEMH}

\subsection{Quy ước mô tả class và method}

Mỗi class cần nêu stereotype, trách nhiệm, thuộc tính quan trọng và invariant. Mỗi method cần nêu signature, mục đích, input/output, precondition, postcondition, exception/error, side effect, transaction boundary và mã requirement/use case liên quan khi phù hợp.

\SubsystemClassDesign{SUB-HUB}{Tra cứu Hub và tài nguyên}
\UseCaseDesign{UC-HUB-01}{Tra cứu Hub phù hợp}{\texttt{SCR-STU-HUB-SEARCH}}{\texttt{FR-HUB-01}, \texttt{BR-HUB-01/02}, \texttt{NIFR-HUB-01}}{xếp hạng theo vị trí; dữ liệu IoT mới/cũ; dịch vụ bản đồ không khả dụng}
\UseCaseDesign{UC-HUB-02}{Xem trạng thái tài nguyên tại Hub}{\texttt{SCR-STU-HUB-DETAIL}}{\texttt{FR-HUB-02}, \texttt{NFR-REL-01}}{snapshot gần nhất, timestamp đồng bộ và tài nguyên \texttt{OutOfService}}
\UseCaseDesign{UC-HUB-03}{Xem Hub thay thế}{\texttt{SCR-STU-HUB-ALT}}{\texttt{FR-HUB-03}, \texttt{EXT-MAP}, \texttt{EXT-STATE}}{không có Hub thay thế; lọc bán kính; dữ liệu trạng thái đến trễ}

\clearpage
\SubsystemClassDesign{SUB-RES}{Đặt phương tiện dùng chung}
\UseCaseDesign{UC-RES-01}{Đặt phương tiện dùng chung}{\texttt{SCR-STU-RES-DETAIL}}{\texttt{FR-RES-01}, \texttt{BR-RES-01--04}, \texttt{NFR-CON-01}}{nhiều request cùng đặt một xe; pin dưới ngưỡng; người dùng đã có reservation}
\UseCaseDesign{UC-RES-02}{Hủy lượt đặt phương tiện}{\texttt{SCR-STU-RES-DETAIL}}{\texttt{FR-RES-02}, \texttt{BR-RES-03}, \texttt{NIFR-RES-01}}{hủy chủ động, hết hạn tự động và giải phóng xe đúng một lần}

\clearpage
\SubsystemClassDesign{SUB-TRIP}{Quản lý chuyến đi}
\UseCaseDesign{UC-TRIP-01}{Nhận phương tiện đã đặt}{\texttt{SCR-STU-TRIP}}{\texttt{FR-TRIP-01}, \texttt{BR-TRIP-01}}{reservation hợp lệ/hết hạn; lệnh mở khóa lỗi; khởi tạo Trip nguyên tử}
\UseCaseDesign{UC-TRIP-02}{Trả phương tiện và kết thúc chuyến đi}{\texttt{SCR-STU-TRIP}}{\texttt{FR-TRIP-02}, \texttt{BR-TRIP-02}}{Hub đích không hợp lệ; khóa/cắm sạc chưa an toàn; chốt số liệu chuyến đi}

\clearpage
\SubsystemClassDesign{SUB-PARK}{Quản lý chỗ đỗ xe cá nhân}
\UseCaseDesign{UC-PARK-01}{Đặt chỗ đỗ cho xe cá nhân}{\texttt{SCR-STU-PARK}}{\texttt{FR-PARK-01}, \texttt{BR-PARK-01/02}, \texttt{NFR-CON-01}}{giữ chỗ đồng thời; thời hạn reservation; Hub đạt ngưỡng quá tải}
\UseCaseDesign{UC-PARK-02}{Check-in phương tiện cá nhân}{\texttt{SCR-STU-PARK}}{\texttt{FR-PARK-02}, \texttt{EXT-STATE}}{cảm biến xác nhận sai chỗ, event lặp và reservation hết hạn}
\UseCaseDesign{UC-PARK-03}{Check-out phương tiện cá nhân}{\texttt{SCR-STU-PARK}}{\texttt{FR-PARK-03}, \texttt{EXT-STATE}}{event rời chỗ đến trễ/lặp và giải phóng chỗ đúng trạng thái}

\clearpage
\SubsystemClassDesign{SUB-CHG}{Quản lý sạc và lịch sạc}
\UseCaseDesign{UC-CHG-01}{Đăng ký nhu cầu sạc}{\texttt{SCR-STU-CHG}, \texttt{SCR-OPR-CHG}}{\texttt{FR-CHG-01}, \texttt{BR-CHG-01/02}}{điểm sạc khả dụng; xếp hàng ưu tiên; request trùng}
\UseCaseDesign{UC-CHG-02}{Theo dõi lịch và phiên sạc}{\texttt{SCR-STU-CHG}, \texttt{SCR-OPR-CHG}}{\texttt{FR-CHG-02}, \texttt{NIFR-CHG-01/02}}{telemetry mới/cũ, hoàn tất phiên và cảnh báo thông số bất thường}
\UseCaseDesign{UC-CHG-03}{Điều chỉnh lịch sạc ưu tiên}{\texttt{SCR-OPR-CHG}}{\texttt{FR-CHG-03}, \texttt{BR-CHG-01/02}}{đổi ưu tiên, tạm dừng/tiếp tục và xung đột lệnh điều phối}

\clearpage
\SubsystemClassDesign{SUB-OPS}{Giám sát và điều phối vận hành}
\UseCaseDesign{UC-OPS-01}{Giám sát mạng lưới Mobility Hub}{\texttt{SCR-OPR-DASHBOARD}}{\texttt{FR-OPS-01}, \texttt{NIFR-OPS-01/02}}{snapshot mạng lưới, ngưỡng mất cân bằng và tài nguyên mất kết nối}
\UseCaseDesign{UC-OPS-02}{Xử lý sự cố vận hành}{\texttt{SCR-OPR-INCIDENT}}{\texttt{FR-OPS-02}, \texttt{BR-OPS-02}}{vòng đời Incident, phân quyền xử lý và audit log}
\UseCaseDesign{UC-OPS-03}{Điều phối phương tiện giữa các Hub}{\texttt{SCR-OPR-REDIST}}{\texttt{FR-OPS-03}, \texttt{BR-OPS-01}}{tạo/phê duyệt kế hoạch, cập nhật tiến độ và số lượng xe không hợp lệ}

\clearpage
\SubsystemClassDesign{SUB-SIM}{What-if Simulation và khuyến nghị}
\UseCaseDesign{UC-SIM-01}{Chạy kịch bản What-if Simulation}{\texttt{SCR-OPR-SIM-CONFIG}}{\texttt{FR-SIM-01}, \texttt{BR-SIM-01}}{validate tham số; snapshot cô lập; lần chạy thất bại hoặc bị hủy}
\UseCaseDesign{UC-SIM-02}{Phân tích kết quả và khuyến nghị điều phối}{\texttt{SCR-OPR-SIM-RESULT}}{\texttt{FR-SIM-02}, \texttt{NIFR-SIM-01}}{sinh khuyến nghị, dữ liệu đầu ra thiếu và không làm thay đổi vận hành thật}

\subsection{Kiểm tra tính nhất quán}

Đối chiếu lifeline trong sequence diagram với các boundary, service, repository và adapter trong class diagram. Transition của state chart phải dẫn tới method thay đổi trạng thái tương ứng. Mỗi UC ở trên phải có ít nhất một đường truy vết hoàn chỉnh từ Screen ID đến sequence lifeline, method xử lý, repository/adapter và test case; không để method nghiệp vụ quan trọng chỉ xuất hiện trong code nhưng vắng mặt trong thiết kế.
